% =================================================================
%  CSE 815 -- Machine Learning
%  Module 2, Week 5, Part 2
%  Multiclass Classification, Softmax, Adam, Convolutional Layers,
%  and Back-propagation
% =================================================================
\documentclass[aspectratio=169,10pt]{beamer}
\usepackage{cse815}

\title[CSE 815 -- Week 5, Part 2]{Multiclass Classification and Back-propagation}
\subtitle{Softmax, a Stable Implementation, Adam, Convolutional Layers, and Computation Graphs}
\author{Rokan Uddin Faruqui\\ \small Professor, Department of CSE}
\institute{\small University of Chittagong \\ \texttt{rufaruqui@cu.ac.bd}}
\date[Week 5 / Part 2]{Week 5 $\cdot$ Part 2}

\begin{document}

\begin{frame}[plain,noframenumbering]
  \titlepage
\end{frame}

% ---------------------------------------------------------------
\begin{frame}{Beyond yes or no}
  \begin{columns}[T,onlytextwidth]
    \begin{column}{0.5\textwidth}
      \begin{block}{Part 1 gave us}
        \footnotesize
        \begin{itemize}
          \item Training: \texttt{Sequential}, \texttt{compile(loss=...)}, \texttt{fit}.
          \item ReLU in hidden layers; output activation chosen from the range of $y$.
          \item Without a non-linear $g$, depth collapses.
        \end{itemize}
      \end{block}
      \footnotesize
      Every output so far was one number. Today the output becomes a probability for
      \emph{each} of $N$ classes. Then we open the two black boxes inside \texttt{fit}: the
      optimiser, and the computation of gradients.
    \end{column}
    \begin{column}{0.46\textwidth}
      \begin{keyidea}[Today]
        \footnotesize
        \begin{enumerate}
          \item Multiclass classification and softmax
          \item Softmax in a network: a stable implementation; multi-label output
          \item Advanced optimisation: Adam
          \item Additional layer types: convolution
          \item Derivatives, computation graphs and back-propagation
        \end{enumerate}
      \end{keyidea}
      \footnotesize The last section is marked optional in the syllabus but is examinable
      here: it is how \texttt{fit} obtains every gradient.
    \end{column}
  \end{columns}
\end{frame}

% ===============================================================
\section{Multiclass Classification and Softmax}

\begin{frame}{More than two classes}
  \begin{columns}[T,onlytextwidth]
    \begin{column}{0.48\textwidth}
      \centering {\scriptsize Illustrative data: four classes in two features}\par
      \begin{tikzpicture}
        \begin{axis}[width=6.4cm,height=4.8cm,xmin=0,xmax=10,ymin=0,ymax=10,
          xlabel={$x_1$},ylabel={$x_2$},xtick=\empty,ytick=\empty,label style={font=\scriptsize}]
          \draw[cugray!60,dashed,thick] (axis cs:5.2,0) -- (axis cs:5.2,10);
          \draw[cugray!60,dashed,thick] (axis cs:0,5.0) -- (axis cs:10,5.6);
          \addplot[only marks,mark=*,mark size=1.8pt,cublue]
            coordinates {(1.5,7.5)(2.4,8.6)(3.1,7.0)(1.9,6.3)(3.6,8.9)(2.8,6.0)};
          \addplot[only marks,mark=x,mark size=2.4pt,thick,curose]
            coordinates {(7.0,8.0)(8.3,7.2)(6.6,6.6)(8.8,8.8)(7.6,9.1)(6.2,7.6)};
          \addplot[only marks,mark=triangle*,mark size=2.2pt,cuteal]
            coordinates {(1.6,2.0)(2.8,3.1)(3.7,1.5)(1.2,3.6)(2.2,1.0)(4.0,2.9)};
          \addplot[only marks,mark=square*,mark size=1.7pt,cuamber]
            coordinates {(7.2,2.2)(8.6,3.4)(6.4,1.4)(9.0,1.6)(7.9,4.1)(6.6,3.2)};
        \end{axis}
      \end{tikzpicture}
    \end{column}
    \begin{column}{0.48\textwidth}
      \footnotesize
      \begin{defnbox}[Multiclass classification]
        \small The target $y$ takes one of $N>2$ \emph{discrete} values: $y\in\{1,\dots,N\}$.
        Exactly one is correct for each example.
      \end{defnbox}
      \begin{block}{Examples}
        \begin{itemize}
          \item Handwritten digit: $N=10$.
          \item Diagnosis of a leaf: healthy, blast, brown spot, or blight.
        \end{itemize}
      \end{block}
      \begin{keyidea}
        We now want an estimate of $P(y=j\mid\vx)$ for every class $j$ --- $N$ numbers
        that are non-negative and sum to $1$.
      \end{keyidea}
    \end{column}
  \end{columns}
\end{frame}

\begin{frame}{The softmax function}
  \begin{columns}[T,onlytextwidth]
    \begin{column}{0.5\textwidth}
      \begin{defnbox}[Softmax regression, $N$ classes]
        \small
        One linear score per class,
        \[ z_j=\vw_j\cdot\vx+b_j,\qquad j=1,\dots,N, \]
        turned into probabilities by
        \[ a_j=\frac{e^{z_j}}{\sum_{k=1}^{N}e^{z_k}}=P(y=j\mid\vx). \]
      \end{defnbox}
      \footnotesize
      Exponentiation makes every score positive; dividing by the sum makes them add to $1$.
      The largest $z_j$ gets the largest $a_j$ --- hence ``soft'' max.
    \end{column}
    \begin{column}{0.46\textwidth}
      \footnotesize
      \begin{block}{Worked example: $\vz=[2,\ 1,\ 0.1]$}
        \renewcommand{\arraystretch}{1.15}
        \begin{tabular}{@{}cccc@{}}
          \toprule
          $j$ & $z_j$ & $e^{z_j}$ & $a_j$\\
          \midrule
          1 & $2$   & $7.3891$ & $0.6590$\\
          2 & $1$   & $2.7183$ & $0.2424$\\
          3 & $0.1$ & $1.1052$ & $0.0986$\\
          \midrule
            &       & $11.2125$ & $1.0000$\\
          \bottomrule
        \end{tabular}
      \end{block}
      \begin{keyidea}[$N=2$ gives back logistic regression]
        \[ a_1=\frac{e^{z_1}}{e^{z_1}+e^{z_2}}=\frac{1}{1+e^{-(z_1-z_2)}}=g(z_1-z_2) \]
        e.g.\ $\vz=[2,0.5]$: $a_1=0.8176=g(1.5)$.
      \end{keyidea}
    \end{column}
  \end{columns}
\end{frame}

\begin{frame}{The softmax loss: cross-entropy}
  \begin{columns}[T,onlytextwidth]
    \begin{column}{0.5\textwidth}
      \begin{defnbox}[Cross-entropy loss]
        \small
        \[ L(\va,y) = -\log a_y =
           \begin{cases}-\log a_1 & y=1\\ \ \ \vdots\\ -\log a_N & y=N\end{cases} \]
        Only the probability assigned to the \emph{correct} class is penalised.
      \end{defnbox}
      \footnotesize
      For $N=2$ this is exactly binary cross-entropy: if $y=1$ the loss is $-\log a_1$; if
      $y=2$ it is $-\log a_2=-\log(1-a_1)$.
    \end{column}
    \begin{column}{0.46\textwidth}
      \centering {\scriptsize $-\log a_y$ against $a_y$}\par
      \begin{tikzpicture}
        \begin{axis}[width=5.8cm,height=3.6cm,xmin=0,xmax=1,ymin=0,ymax=4.5,
          xlabel={$a_y$},tick label style={font=\scriptsize},label style={font=\scriptsize}]
          \addplot[curose,very thick,domain=0.012:1,samples=120]{-ln(x)};
          \addplot[only marks,mark=*,mark size=1.6pt,cublue] coordinates
            {(0.6590,0.4170)(0.2424,1.4170)(0.0986,2.3170)};
        \end{axis}
      \end{tikzpicture}
      \footnotesize
      \begin{block}{Same example}
        $L=0.4170$ if $y=1$; \ $1.4170$ if $y=2$; \ $2.3170$ if $y=3$.\par
        The losses differ by exactly the gaps between the $z_j$: \
        $-\log a_j=\log\sum_k e^{z_k}-z_j$.
      \end{block}
    \end{column}
  \end{columns}
\end{frame}

% ===============================================================
\section{Softmax in a Neural Network}

\begin{frame}[fragile]{A softmax output layer}
  \begin{columns}[T,onlytextwidth]
    \begin{column}{0.48\textwidth}
      \centering
      \begin{tikzpicture}[node distance=3mm,
          lay/.style={draw=cublue,fill=culight,rounded corners,font=\tiny,
                      minimum height=5mm,inner sep=2pt}]
        \node[lay,draw=cuteal,fill=cuteal!12] (x) {$\vx$: 784};
        \node[lay,right=of x] (l1) {25 ReLU};
        \node[lay,right=of l1] (l2) {15 ReLU};
        \node[lay,right=of l2,draw=cuamber,fill=cuamber!15] (l3) {10 softmax};
        \foreach \a/\b in {x/l1,l1/l2,l2/l3}{\draw[synapse] (\a)--(\b);}
      \end{tikzpicture}\par
      \vspace{2mm}
      \footnotesize\raggedright
      Digits $0$--$9$ from $28\times28$ images. The output layer computes
      \[ z^{[3]}_j=\vw^{[3]}_j\cdot\va^{[2]}+b^{[3]}_j,\qquad
         a^{[3]}_j=\frac{e^{z^{[3]}_j}}{\sum_{k=1}^{10}e^{z^{[3]}_k}} \]
\begin{lstlisting}
from tensorflow.keras.losses import \
    SparseCategoricalCrossentropy
model = Sequential([
    Dense(units=25, activation='relu'),
    Dense(units=15, activation='relu'),
    Dense(units=10, activation='softmax')])
model.compile(loss=SparseCategoricalCrossentropy())
model.fit(X, y, epochs=100)
\end{lstlisting}
    \end{column}
    \begin{column}{0.48\textwidth}
      \footnotesize
      \begin{keyidea}[Softmax is different]
        Sigmoid, ReLU and linear act on each $z_j$ alone. Each softmax output $a_j$ depends
        on \emph{all} of $z_1,\dots,z_N$, through the denominator. Raise one score and every
        other probability falls.
      \end{keyidea}
      \begin{defnbox}[``Sparse categorical'']
        \small \emph{Categorical}: $y$ is one of $N$ categories. \emph{Sparse}: $y$ is
        stored as a single integer ($0$--$9$), not as a length-$10$ one-hot vector.
      \end{defnbox}
      \begin{pitfall}[Do not ship this version]
        \footnotesize It works, but it is numerically fragile. The next two slides give
        the version you should actually use.
      \end{pitfall}
    \end{column}
  \end{columns}
\end{frame}

\begin{frame}[fragile]{Round-off: the same number, two ways}
  \begin{columns}[T,onlytextwidth]
    \begin{column}{0.5\textwidth}
\begin{lstlisting}
x1 = 2.0 / 10000
print(f"{x1:.20f}")
# 0.00020000000000000001

x2 = 1 + 1/10000 - (1 - 1/10000)
print(f"{x2:.20f}")
# 0.00019999999999997797
\end{lstlisting}
      \footnotesize
      Mathematically $x_1=x_2$. In floating point, computing a small number as the
      difference of two numbers near $1$ loses digits.
      \begin{block}{The softmax version of the problem}
        For logits $\vz=[1000,999,998]$, \texttt{np.exp(1000)} overflows to \texttt{inf},
        and $a_1=\texttt{inf}/\texttt{inf}=\texttt{nan}$; very negative $z$ give $\log 0$.
      \end{block}
    \end{column}
    \begin{column}{0.46\textwidth}
      \footnotesize
      \begin{keyidea}[Do not compute $a$ and then $\log a$]
        Substitute the softmax into the loss and simplify \emph{before} evaluating:
        \[ -\log a_y = \log\sum_{k}e^{z_k}-z_y. \]
        Shift by $c=\max_k z_k$, so that every exponent is $\le 0$:
        \[ \log\sum_k e^{z_k} = c+\log\sum_k e^{z_k-c}. \]
      \end{keyidea}
      \begin{block}{Check it}
        $\vz=[1000,999,998]$, $y=1$:
        \[ L=\log(1+e^{-1}+e^{-2})=0.4076 \]
        --- finite, and computed without ever forming $e^{1000}$.
      \end{block}
    \end{column}
  \end{columns}
\end{frame}

\begin{frame}[fragile]{The recommended implementation: \texttt{from\_logits=True}}
  \begin{columns}[T,onlytextwidth]
    \begin{column}{0.52\textwidth}
\begin{lstlisting}
model = Sequential([
    Dense(units=25, activation='relu'),
    Dense(units=15, activation='relu'),
    Dense(units=10, activation='linear')])  # z, not a
model.compile(
    loss=SparseCategoricalCrossentropy(from_logits=True))
model.fit(X, y, epochs=100)

logits = model(X_new)            # z: scores, not probs
probs  = tf.nn.softmax(logits)   # if probabilities needed
yhat   = np.argmax(logits, axis=1)
\end{lstlisting}
      \vspace{1mm}
      \footnotesize\textbf{Binary case, same idea:} \texttt{Dense(1, 'linear')} with
      \texttt{BinaryCrossentropy(from\_logits=True)}; recover the probability with
      \texttt{tf.nn.sigmoid(logits)}.
    \end{column}
    \begin{column}{0.44\textwidth}
      \footnotesize
      \begin{defnbox}[Logits]
        \small The values $\vz$ that the last layer outputs \emph{before} the softmax (or
        sigmoid). \texttt{from\_logits=True} tells the loss to apply the softmax itself,
        using the stable formula.
      \end{defnbox}
      \begin{pitfall}[The model now outputs $\vz$]
        \footnotesize \texttt{model(X)} returns logits, which can be negative and do not
        sum to $1$. Apply \texttt{tf.nn.softmax} before reporting probabilities. For the
        predicted class you do not need it: $\argmax_j z_j=\argmax_j a_j$.
      \end{pitfall}
    \end{column}
  \end{columns}
\end{frame}

\begin{frame}{Multi-label classification (several outputs)}
  \begin{columns}[T,onlytextwidth]
    \begin{column}{0.5\textwidth}
      \begin{block}{A different question}
        \footnotesize
        A street camera image. Is there a \textbf{car}? A \textbf{bus}? A
        \textbf{pedestrian}? Any combination may be present, so the label is a vector:
        \[ \vec{y}=[1,\ 0,\ 1] \quad\text{(car, no bus, pedestrian)}. \]
      \end{block}
      \centering
      \begin{tikzpicture}[node distance=3mm,
          lay/.style={draw=cublue,fill=culight,rounded corners,font=\tiny,
                      minimum height=5mm,inner sep=2pt}]
        \node[lay,draw=cuteal,fill=cuteal!12] (x) {image $\vx$};
        \node[lay,right=of x] (l1) {hidden};
        \node[lay,right=of l1] (l2) {hidden};
        \node[lay,right=of l2,draw=cuamber,fill=cuamber!15,align=center] (l3)
          {3 sigmoid\\[-1pt]units};
        \foreach \a/\b in {x/l1,l1/l2,l2/l3}{\draw[synapse] (\a)--(\b);}
      \end{tikzpicture}
    \end{column}
    \begin{column}{0.46\textwidth}
      \footnotesize
      \renewcommand{\arraystretch}{1.2}
      \begin{tabular}{@{}L{1.6cm}L{2.2cm}L{2.4cm}@{}}
        \toprule
        & multiclass & multi-label\\
        \midrule
        labels per example & exactly one & any number\\
        output layer & $N$ softmax units & $N$ sigmoid units\\
        outputs sum to & $1$ & anything in $[0,N]$\\
        loss & cross-entropy & binary cross-entropy per output, summed\\
        \bottomrule
      \end{tabular}
      \begin{pitfall}[Choosing by the data, not by habit]
        \footnotesize Softmax on a multi-label task forces the classes to compete: a
        confident ``car'' pushes ``pedestrian'' down even when both are present.
      \end{pitfall}
    \end{column}
  \end{columns}
\end{frame}

% ===============================================================
\section{Advanced Optimisation: Adam}

\begin{frame}{One learning rate does not suit every parameter}
  \begin{columns}[T,onlytextwidth]
    \begin{column}{0.5\textwidth}
      \centering {\scriptsize $J=w_1^2+10\,w_2^2$, gradient descent from the left}\par
      \begin{tikzpicture}
        \begin{axis}[name=lo,width=4.4cm,height=3.7cm,view={0}{90},
          xmin=-4.5,xmax=1,ymin=-1.3,ymax=1.3,title={\tiny $\alpha=0.01$: creeping},
          tick label style={font=\tiny},xtick={-4,-2,0},ytick={-1,0,1}]
          \costcontour{cugray!60,forget plot,domain=-4.5:1,y domain=-1.3:1.3}{1,4,9,16}{x^2+10*y^2}
          \addplot[cuteal,thick,mark=*,mark size=0.9pt] coordinates
            {(-4.000,0.800)(-3.920,0.640)(-3.842,0.512)(-3.765,0.410)(-3.689,0.328)
             (-3.616,0.262)(-3.543,0.210)(-3.473,0.168)(-3.403,0.134)(-3.335,0.107)
             (-3.268,0.086)(-3.203,0.069)(-3.139,0.055)(-3.076,0.044)(-3.015,0.035)};
        \end{axis}
        \begin{axis}[at={(lo.east)},anchor=west,xshift=4mm,width=4.4cm,height=3.7cm,
          view={0}{90},xmin=-4.5,xmax=1,ymin=-1.3,ymax=1.3,
          title={\tiny $\alpha=0.09$: zig-zag},
          tick label style={font=\tiny},xtick={-4,-2,0},ytick={-1,0,1},yticklabels={}]
          \costcontour{cugray!60,forget plot,domain=-4.5:1,y domain=-1.3:1.3}{1,4,9,16}{x^2+10*y^2}
          \addplot[curose,thick,mark=*,mark size=0.9pt] coordinates
            {(-4.000,1.000)(-3.280,-0.800)(-2.690,0.640)(-2.205,-0.512)(-1.808,0.410)
             (-1.483,-0.328)(-1.216,0.262)(-0.997,-0.210)(-0.818,0.168)(-0.670,-0.134)
             (-0.550,0.107)(-0.451,-0.086)};
        \end{axis}
      \end{tikzpicture}\par
      {\scriptsize $x$-axis $w_1$, $y$-axis $w_2$; 15 and 12 steps shown.}
    \end{column}
    \begin{column}{0.46\textwidth}
      \footnotesize
      \begin{block}{Two failure modes, one fix}
        \textbf{Left:} every step moves $w_1$ in the same direction, but slowly --- $w_1$
        wants a \emph{larger} step.\\[2pt]
        \textbf{Right:} $w_2$ flips sign every step --- it wants a \emph{smaller} step,
        while $w_1$ is still happy to go faster.
      \end{block}
      \begin{keyidea}[Adam --- adaptive moment estimation]
        Give every parameter its own effective step size, adjusted automatically from the
        history of its gradients:
        \begin{itemize}
          \item gradient keeps the same sign $\Rightarrow$ steps stay large;
          \item gradient oscillates $\Rightarrow$ the average cancels, steps shrink.
        \end{itemize}
      \end{keyidea}
    \end{column}
  \end{columns}
\end{frame}

\begin{frame}[fragile]{The Adam update, and how to use it}
  \begin{columns}[T,onlytextwidth]
    \begin{column}{0.52\textwidth}
      \begin{defnbox}[Adam, for one parameter $w$ at step $t$]
        \small
        With gradient $d=\partial J/\partial w$:
        \begin{align*}
          m &\leftarrow \beta_1 m+(1-\beta_1)\,d && \text{mean of } d\\
          v &\leftarrow \beta_2 v+(1-\beta_2)\,d^2 && \text{mean of } d^2\\
          w &\leftarrow w-\alpha\,\frac{\hat m}{\sqrt{\hat v}+\epsilon},
          && \hat m=\tfrac{m}{1-\beta_1^t},\ \hat v=\tfrac{v}{1-\beta_2^t}
        \end{align*}
        Defaults: $\beta_1=0.9$, $\beta_2=0.999$, $\epsilon=10^{-7}$; $\alpha=10^{-3}$.
      \end{defnbox}
      \footnotesize
      $\hat m/\sqrt{\hat v}$ is near $\pm1$ when the gradient is consistent, so each step is
      about $\alpha$ \emph{whatever the gradient's scale}; near $0$ when it oscillates.
    \end{column}
    \begin{column}{0.44\textwidth}
      \footnotesize
\begin{lstlisting}
model.compile(
  optimizer=tf.keras.optimizers.Adam(
      learning_rate=1e-3),
  loss=SparseCategoricalCrossentropy(
      from_logits=True))
\end{lstlisting}
      \vspace{1mm}
      \begin{block}{In practice}
        Adam is the default optimiser for neural networks. It is far less sensitive to the
        choice of $\alpha$ than plain gradient descent, but $\alpha$ still matters: try
        $10^{-2}$, $10^{-3}$, $10^{-4}$.
      \end{block}
      \begin{pitfall}[Not magic]
        \footnotesize A bad loss, unscaled features or a bug in the data are not fixed by
        changing the optimiser. Check the learning curve first.
      \end{pitfall}
    \end{column}
  \end{columns}
\end{frame}

% ===============================================================
\section{Additional Layer Types}

\begin{frame}{Dense layers, and an alternative}
  \begin{columns}[T,onlytextwidth]
    \begin{column}{0.5\textwidth}
      \begin{block}{Dense layer}
        \footnotesize Every unit sees every activation of the previous layer. Powerful,
        but parameter-hungry: $784$ pixels into $25$ units already costs $19\,625$
        parameters.
      \end{block}
      \begin{defnbox}[Convolutional layer]
        \small Each unit looks only at a \textbf{window} --- a small, contiguous region of
        its input --- and the \emph{same} weights (a \textbf{filter}) are slid across all
        windows.
      \end{defnbox}
      \footnotesize
      \textbf{Why it helps.}
      \begin{itemize}
        \item Far fewer parameters $\Rightarrow$ faster to compute.
        \item Needs less data, and is less prone to overfitting (Week 3).
        \item Builds in a useful assumption: a pattern means the same thing wherever it
              appears.
      \end{itemize}
    \end{column}
    \begin{column}{0.46\textwidth}
      \centering {\scriptsize Each hidden unit sees one region of the image}\par
      \begin{tikzpicture}[scale=0.3]
        \draw[cugray!60] (0,0) grid (8,8);
        \fill[cuteal,opacity=0.25] (0,4) rectangle (4,8);
        \fill[cuamber,opacity=0.25] (4,4) rectangle (8,8);
        \fill[curose,opacity=0.25] (2,0) rectangle (6,4);
        \node[neuron,minimum size=5mm] (u1) at (12,7) {};
        \node[neuron,minimum size=5mm] (u2) at (12,4.5) {};
        \node[neuron,minimum size=5mm] (u3) at (12,2) {};
        \draw[synapse,cuteal] (2,6) -- (u1);
        \draw[synapse,cuamber] (6,6) -- (u2);
        \draw[synapse,curose] (4,2) -- (u3);
      \end{tikzpicture}
      \vspace{3mm}

      \footnotesize\raggedright
      Hand-designing the window size, the stride (how far each window moves), and the
      number of filters is part of the architecture. Convolutional networks are the
      workhorse of computer vision.
    \end{column}
  \end{columns}
\end{frame}

\begin{frame}{Example: classifying a heart signal}
  \begin{columns}[T,onlytextwidth]
    \begin{column}{0.54\textwidth}
      \centering {\scriptsize An EKG trace: $100$ samples $x_1,\dots,x_{100}$ (schematic)}\par
      \begin{tikzpicture}
        \begin{axis}[width=7.4cm,height=2.9cm,xmin=0,xmax=100,ymin=-1.2,ymax=2.3,clip=false,
          axis lines=none]
          \addplot[cublue,thick,domain=0:100,samples=400]
            {0.25*sin(deg(x/3))+1.6*exp(-((mod(x,25)-12)^2)/1.2)-0.4*exp(-((mod(x,25)-15)^2)/2)};
          \draw[cuteal,thick,decorate,decoration={brace,mirror,amplitude=3pt}]
            (axis cs:0,-0.9) -- node[below=2pt,font=\tiny]{unit 1: $x_1..x_{20}$} (axis cs:20,-0.9);
          \draw[cuamber,thick,decorate,decoration={brace,amplitude=3pt}]
            (axis cs:10,2.05) -- node[above=2pt,font=\tiny]{unit 2: $x_{11}..x_{30}$} (axis cs:30,2.05);
          \draw[curose,thick,decorate,decoration={brace,mirror,amplitude=3pt}]
            (axis cs:80,-0.9) -- node[below=2pt,font=\tiny]{unit 9: $x_{81}..x_{100}$} (axis cs:100,-0.9);
        \end{axis}
      \end{tikzpicture}\par
      \vspace{4mm}
      \begin{tikzpicture}[node distance=3mm,
          lay/.style={draw=cublue,fill=culight,rounded corners,font=\tiny,
                      minimum height=5mm,inner sep=2pt,align=center}]
        \node[lay,draw=cuteal,fill=cuteal!12] (x) {$\vx$\\100};
        \node[lay,right=of x] (l1) {conv: 9 units\\window 20, stride 10};
        \node[lay,right=of l1] (l2) {conv: 3 units\\window 5, stride 2};
        \node[lay,right=of l2,draw=cuamber,fill=cuamber!15] (l3) {1 sigmoid\\heart disease?};
        \foreach \a/\b in {x/l1,l1/l2,l2/l3}{\draw[synapse] (\a)--(\b);}
      \end{tikzpicture}
    \end{column}
    \begin{column}{0.42\textwidth}
      \footnotesize
      \begin{block}{Windows}
        Layer 1: windows start at $1,11,\dots,81$: $9$ units. Layer 2 looks at
        $a^{[1]}_{1..5}$, $a^{[1]}_{3..7}$, $a^{[1]}_{5..9}$: $3$ units.
      \end{block}
      \begin{block}{Parameters in layer 1}
        \renewcommand{\arraystretch}{1.15}
        \begin{tabular}{@{}ll@{}}
          dense, $100\to9$ & $100\cdot9+9=909$\\
          own weights per window & $9\cdot(20+1)=189$\\
          one shared filter & $20+1=21$\\
        \end{tabular}
      \end{block}
      \begin{keyidea}
        A real convolutional layer uses several filters, each shared across all windows,
        so it detects several patterns anywhere in the signal for a few dozen parameters
        apiece.
      \end{keyidea}
    \end{column}
  \end{columns}
\end{frame}

% ===============================================================
\section{Derivatives, Computation Graphs and Back-propagation}

\begin{frame}{What a derivative tells gradient descent}
  \begin{columns}[T,onlytextwidth]
    \begin{column}{0.5\textwidth}
      \begin{block}{Informal definition}
        \footnotesize
        If nudging $w$ by a tiny $\varepsilon$ changes $J(w)$ by about $k\varepsilon$, then
        \[ \frac{dJ}{dw}=k. \]
        Example: $J=w^2$ at $w=3$, with $\varepsilon=0.001$:\par
        $J(3.001)=9.006001$, a rise of $0.006001\approx 6\varepsilon$. So $dJ/dw=6$ ---
        and indeed $\tfrac{d}{dw}w^2=2w=6$.
      \end{block}
      \footnotesize
      The update $w\leftarrow w-\alpha\,dJ/dw$: a large derivative means $J$ is sensitive
      to $w$, so take a large step; a negative one means increasing $w$ lowers $J$.
    \end{column}
    \begin{column}{0.46\textwidth}
      \footnotesize
      \begin{block}{Check: derivatives at $w=2$, $\varepsilon=0.001$}
        \renewcommand{\arraystretch}{1.2}
        \begin{tabular}{@{}llll@{}}
          \toprule
          $J(w)$ & $\dfrac{dJ}{dw}$ & at $w=2$ & $\dfrac{J(2.001)-J(2)}{0.001}$\\
          \midrule
          $w^2$   & $2w$        & $4$     & $4.001$\\
          $w^3$   & $3w^2$      & $12$    & $12.006$\\
          $w$     & $1$         & $1$     & $1$\\
          $1/w$   & $-1/w^2$    & $-0.25$ & $-0.2499$\\
          \bottomrule
        \end{tabular}
      \end{block}
      \begin{keyidea}[Numerical gradient check]
        Comparing an analytic derivative with $\bigl(J(w+\varepsilon)-J(w)\bigr)/\varepsilon$
        is the standard way to test gradient code. Symbolic tools (e.g.\ SymPy) do the
        algebra if you need it.
      \end{keyidea}
    \end{column}
  \end{columns}
\end{frame}

\begin{frame}{The computation graph}
  \footnotesize
  \centering
  Linear regression on one example: $J=\tfrac12(wx+b-y)^2$, with $w=2,\ b=8,\ x=-2,\ y=2$.\par
  \vspace{2mm}
  \begin{tikzpicture}[
      op/.style={draw=cublue,fill=culight,rounded corners,minimum width=1.9cm,
                 minimum height=7mm,align=center,font=\scriptsize},
      inp/.style={font=\scriptsize,cuteal}]
    \node[inp] (w) at (0,0.55) {$w=2$};
    \node[inp] (x) at (0,-0.55) {$x=-2$};
    \node[op] (c) at (2.5,0) {$c=wx$\\$=-4$};
    \node[inp] (b) at (2.5,-1.35) {$b=8$};
    \node[op] (a) at (5.3,0) {$a=c+b$\\$=4$};
    \node[inp] (y) at (5.3,-1.35) {$y=2$};
    \node[op] (d) at (8.1,0) {$d=a-y$\\$=2$};
    \node[op,draw=cuamber,fill=cuamber!12] (J) at (10.9,0) {$J=\tfrac12d^2$\\$=2$};
    \draw[synapse,thick] (w) -- (c); \draw[synapse,thick] (x) -- (c);
    \draw[synapse,thick] (c) -- (a); \draw[synapse,thick] (b) -- (a);
    \draw[synapse,thick] (a) -- (d); \draw[synapse,thick] (y) -- (d);
    \draw[synapse,thick] (d) -- (J);
    \foreach \s/\t/\l in {J/d/{$\frac{\partial J}{\partial d}=d=2$},
                          d/a/{$\frac{\partial J}{\partial a}=2$},
                          a/c/{$\frac{\partial J}{\partial c}=2$}}{
      \draw[-{Stealth},curose,thick] ([yshift=4mm]\s.north) to[bend right=30]
        node[above,font=\tiny,curose]{\l} ([yshift=4mm]\t.north);}
    \node[font=\tiny,curose,align=left] at (1.0,1.25) {$\frac{\partial J}{\partial w}=\frac{\partial J}{\partial c}\cdot x=-4$};
    \node[font=\tiny,curose] at (3.9,-1.35) {$\frac{\partial J}{\partial b}=\frac{\partial J}{\partial a}=2$};
  \end{tikzpicture}\par
  \vspace{2mm}
  \raggedright
  \begin{columns}[T,onlytextwidth]
    \begin{column}{0.5\textwidth}
      \begin{block}{Forward, left to right}
        Compute and store every intermediate value: $c$, $a$, $d$, $J$.
      \end{block}
    \end{column}
    \begin{column}{0.46\textwidth}
      \begin{block}{Backward, right to left}
        Each node multiplies the derivative arriving from its right by its own local
        derivative (chain rule). Check: nudging $w$ to $2.001$ changes $J$ by $-0.004$.
      \end{block}
    \end{column}
  \end{columns}
\end{frame}

\begin{frame}{Why backward is efficient}
  \begin{columns}[T,onlytextwidth]
    \begin{column}{0.5\textwidth}
      \begin{keyidea}[Reuse]
        $\partial J/\partial a$ was computed once and then used for \emph{both}
        $\partial J/\partial b$ and $\partial J/\partial c$ (and through $c$, for $w$).
        Every intermediate derivative is shared by everything to its left.
      \end{keyidea}
      \footnotesize
      \begin{block}{Cost, for $N$ nodes and $P$ parameters}
        \begin{itemize}
          \item Nudge each parameter and re-run the forward pass: about $N\times P$ steps.
          \item Back-propagation: about $N+P$ steps.
        \end{itemize}
        With $N=10^4$ and $P=10^5$: $10^9$ against $1.1\times10^5$ --- the difference
        between impossible and instant.
      \end{block}
    \end{column}
    \begin{column}{0.46\textwidth}
      \footnotesize
      \begin{defnbox}[Back-propagation]
        \small Compute all $\partial J/\partial(\text{parameter})$ in one backward sweep over
        the computation graph, applying the chain rule at each node and reusing every
        intermediate result.
      \end{defnbox}
      \begin{block}{Automatic differentiation}
        TensorFlow records the forward computation as a graph while it runs, then runs
        back-propagation over it. This is the ``library does the calculus'' of Part 1 ---
        and why you only ever write the forward pass.
      \end{block}
    \end{column}
  \end{columns}
\end{frame}

\begin{frame}{A larger example: two ReLU layers}
  \footnotesize
  \begin{columns}[T,onlytextwidth]
    \begin{column}{0.44\textwidth}
      \begin{block}{Setting}
        One input, one unit per layer, ReLU in both, squared-error loss:
        \[ a^{[1]}=g(w^{[1]}x+b^{[1]}),\quad a^{[2]}=g(w^{[2]}a^{[1]}+b^{[2]}) \]
        \[ J=\tfrac12\bigl(a^{[2]}-y\bigr)^2 \]
        $x=1$, $y=5$; \ $w^{[1]}=2$, $b^{[1]}=0$; \ $w^{[2]}=3$, $b^{[2]}=1$.
      \end{block}
      \begin{block}{Forward}
        $z^{[1]}=2$, $a^{[1]}=2$; \ $z^{[2]}=3\cdot2+1=7$, $a^{[2]}=7$;\\
        $J=\tfrac12(7-5)^2=2$.
      \end{block}
    \end{column}
    \begin{column}{0.52\textwidth}
      \begin{block}{Backward (ReLU$'(z)=1$ here, since every $z>0$)}
        \renewcommand{\arraystretch}{1.25}
        \begin{tabular}{@{}lll@{}}
          $\partial J/\partial a^{[2]}$ & $=a^{[2]}-y$ & $=2$\\
          $\partial J/\partial z^{[2]}$ & $=2\cdot1$ & $=2$\\
          $\partial J/\partial w^{[2]}$ & $=2\cdot a^{[1]}$ & $=\mathbf{4}$\\
          $\partial J/\partial b^{[2]}$ & $=2$ & $=\mathbf{2}$\\
          $\partial J/\partial a^{[1]}$ & $=2\cdot w^{[2]}$ & $=6$\\
          $\partial J/\partial z^{[1]}$ & $=6\cdot1$ & $=6$\\
          $\partial J/\partial w^{[1]}$ & $=6\cdot x$ & $=\mathbf{6}$\\
          $\partial J/\partial b^{[1]}$ & $=6$ & $=\mathbf{6}$\\
        \end{tabular}
      \end{block}
      \begin{pitfall}[Verify numerically]
        \footnotesize Nudging each parameter by $0.001$ gives $6.0045$, $6.0045$,
        $4.0020$, $2.0005$ --- agreeing to the size of $\varepsilon$.
      \end{pitfall}
    \end{column}
  \end{columns}
\end{frame}

% ===============================================================
\section{Wrap-up and Assessment}

\begin{frame}{Summary}
  \begin{columns}[T,onlytextwidth]
    \begin{column}{0.5\textwidth}
      \begin{block}{What to remember}
        \footnotesize
        \begin{itemize}
          \item Softmax turns $N$ scores into $N$ probabilities; with $N=2$ it is logistic
                regression.
          \item Loss $-\log a_y$: \texttt{SparseCategoricalCrossentropy}.
          \item Output linear logits and pass \texttt{from\_logits=True}: numerically stable.
          \item Multi-label: $N$ sigmoid outputs, not softmax.
          \item Adam adapts a step size per parameter; it is the default optimiser.
          \item Convolutional units see windows and share weights.
          \item Back-propagation: one backward sweep, chain rule, $\approx N+P$ steps.
        \end{itemize}
      \end{block}
    \end{column}
    \begin{column}{0.46\textwidth}
      \begin{block}{Formulas}
        \[ a_j=\frac{e^{z_j}}{\sum_k e^{z_k}},\qquad L=-\log a_y \]
        \[ -\log a_y = \log\sum_k e^{z_k}-z_y \]
        \[ \frac{\partial J}{\partial w}=\frac{\partial J}{\partial c}\cdot\frac{\partial c}{\partial w} \]
      \end{block}
      \begin{block}{Next}
        \footnotesize Week 6: advice for applying machine learning --- evaluating a model,
        train/validation/test sets, and diagnosing bias and variance.
      \end{block}
    \end{column}
  \end{columns}
\end{frame}

\begin{frame}{Descriptive questions}
  \footnotesize
  \begin{enumerate}
    \item Define multiclass classification and give two examples. Why can one sigmoid
          output not handle it? \hfill\clo{4}
    \item State the softmax function. Compute $\va$ for $\vz=[1,\ 0,\ -1]$ and verify that
          the outputs sum to $1$. \hfill\clo{6}
    \item Show that softmax with $N=2$ is the same as logistic regression, and say what
          plays the role of $z$. \hfill\clo{6}
    \item Write the cross-entropy loss for softmax. For the worked example
          $\vz=[2,1,0.1]$, which label gives the largest loss, and why? \hfill\clo{5}
    \item Explain why softmax is unlike sigmoid, ReLU and linear activations in how each
          output depends on the scores. \hfill\clo{4}
    \item What are logits? Explain why \texttt{from\_logits=True} with a linear output layer
          is preferred, and how you then obtain probabilities. \hfill\clo{6}
  \end{enumerate}
\end{frame}

\begin{frame}{Descriptive questions (continued)}
  \footnotesize
  \begin{enumerate}\setcounter{enumi}{6}
    \item Distinguish multiclass from multi-label classification, giving the output layer
          and loss for each. \hfill\clo{5}
    \item Describe the Adam update. Explain how it behaves when a parameter's gradient
          keeps the same sign, and when it oscillates. \hfill\clo{6}
    \item What is a convolutional layer? A $1$-D input has $200$ samples; layer 1 uses
          windows of $40$ with stride $20$. How many units, and how many parameters with one
          shared filter? \hfill\clo{6}
    \item For $J=\tfrac12(wx+b-y)^2$ with $w=1$, $b=2$, $x=3$, $y=1$, draw the computation
          graph and compute $\partial J/\partial w$ and $\partial J/\partial b$ by
          back-propagation. \hfill\clo{8}
    \item Explain why back-propagation takes about $N+P$ steps whereas finite differences
          take about $N\times P$. \hfill\clo{5}
    \item In the two-ReLU-layer example, recompute every derivative with $w^{[2]}=-1$.
          What happens, and why? \hfill\clo{8}
  \end{enumerate}
\end{frame}

\begin{frame}{Design questions}
  \footnotesize
  \begin{enumerate}
    \item \textbf{Design a crop-disease classifier.} The Bangladesh Rice Research Institute
          wants to classify leaf photographs as healthy, blast, brown spot or bacterial
          blight. Specify the architecture (including any convolutional layers), every
          activation, the loss, the optimiser, and how you will output probabilities to
          farmers. \hfill\clo{12}
    \item \textbf{Design a news tagger.} A Bangla news portal wants to tag each article with
          any of $12$ topics (politics, sport, economy, \dots). Design the output layer and
          loss, and explain what would go wrong with softmax. \hfill\clo{10}
    \item \textbf{Design a gradient checker.} Design a routine that verifies the gradients
          produced by a back-propagation implementation for an arbitrary network. Specify
          $\varepsilon$, the comparison metric, the tolerance, and how you handle the
          ReLU kink at $0$. \hfill\clo{12}
    \item \textbf{Design an optimiser comparison.} Design an experiment comparing plain
          gradient descent and Adam on the digit data. Specify the learning rates you
          sweep, what you measure, and how you make the comparison fair. \hfill\clo{10}
  \end{enumerate}
\end{frame}

\begin{frame}{Reading}
  \begin{block}{Primary text}
    G.~James, D.~Witten, T.~Hastie and R.~Tibshirani,
    \emph{An Introduction to Statistical Learning}, 2nd ed., Springer, 2021.\\[2pt]
    \hfill{\small$\rightarrow$ \textbf{\S10.2} --- softmax output and cross-entropy on the
    MNIST digits; \textbf{\S10.3} --- convolutional networks; \textbf{\S10.7.1} ---
    back-propagation.}
  \end{block}
  \begin{block}{Supplementary}
    C.~M.~Bishop, \emph{Pattern Recognition and Machine Learning}, Springer, 2006.\\[2pt]
    \hfill{\small$\rightarrow$ \textbf{\S4.3.4} --- multiclass logistic regression
    (softmax); \textbf{\S5.3} --- error back-propagation.}
  \end{block}
  \vspace{1mm}
  \begin{center}
    \small\color{cugray}
    Lab (CSE 816, Week 5): softmax; multiclass; practice lab on neural networks for
    multiclass classification; derivatives and back-propagation.\\
    Practice quizzes: \emph{Multiclass Classification}; \emph{Additional Neural Network
    Concepts}.
  \end{center}
\end{frame}

\end{document}
